\section{Full local Galois lifts and simultaneous minimum layers}
\label{sec:full-galois-lifts}

Proposition~\ref{prop:trace-transport} gives a necessary trace constraint
for local Galois transport. We now construct actual automorphisms of full local absolute
Galois groups and compute their action on specified vectors. The two
examples below give a common minimum layer for a background vector and,
in the second example, three different square-index vectors. All tensor
factors use the same local arrow. The resulting exact maximal-order
hulls concern the native families defined here; their identification
with a published global collation is a separate question.

Throughout this section, $p$ is odd, $E/\Q_p$ is finite, and the native
valuation is normalized by $v_p(p)=1$. Write
\[
 G_E=\operatorname{Gal}(\overline E/E),\qquad
 L_E=\log\OO_E^\times,\qquad I_E=p^{-1}L_E.
\]
The notation for free generators and local coordinate vectors in this
section is independent of the global endpoints of an abc triple.

\subsection{A word automorphism and the full relative presentation}

Let $n=[E:\Q_p]$ be even. The original Jannsen--Wingberg construction
\cite[\S1.2, Theorem~2, and \S1.4(a), pp.~74--76]{JannsenWingberg}
uses the full tame group $\mathfrak T$, generated by $\sigma,\theta$,
and the relative profinite group
\begin{equation}\label{eq:gl-relative}
 V=\widehat F_{n+1}*\mathfrak T,\qquad
 Z=\ker(V\longrightarrow\mathfrak T),\qquad
 J=\ker(Z\longrightarrow Z^{(p)}),\qquad
 \mathcal F=V/J.
\end{equation}
Here $\widehat F_{n+1}$ is free profinite on $x_0,\ldots,x_n$, and
$Z^{(p)}$ is the maximal pro-$p$ quotient of $Z$. One then imposes
the additional relation
\begin{equation}\label{eq:gl-jw}
 \sigma x_0\sigma^{-1}
 =W_0(x_0,\theta)x_1^{p^s}
    [x_1,x_2][x_3,x_4]\cdots[x_{n-1},x_n],
\end{equation}
where $W_0$ belongs to the closed subgroup generated by $x_0,\theta$.
The commutator convention is $[x,y]=xyx^{-1}y^{-1}$.
The tame relation $\sigma\theta\sigma^{-1}=\theta^{p^f}$ and all
internal relations of $\mathfrak T$ are retained in
\eqref{eq:gl-relative}. In particular, this construction is not the
maximal pro-$p$ quotient of the whole absolute Galois group.

\begin{lemma}[An exact cross-handle word]\label{lem:gl-word}
In the free group on $a,b,c,d_0$, put
\[
 r=bab^{-1},\qquad z=d_0rd_0^{-1},\qquad
 \delta=[a,b][c,d_0].
\]
The substitution
\begin{equation}\label{eq:gl-forward}
 \begin{aligned}
 F(a)&=a,& F(b)&=d_0b,\\
 F(c)&=zr^{-1}cz^{-1},&F(d_0)&=zd_0z^{-1}
 \end{aligned}
\end{equation}
is an automorphism fixing $a$ and $\delta$ exactly. Its inverse is
\begin{equation}\label{eq:gl-inverse}
 \begin{aligned}
 G(a)&=a,&G(b)&=r^{-1}d_0^{-1}rb,\\
 G(c)&=r^{-1}d_0^{-1}rd_0cr,&G(d_0)&=r^{-1}d_0r.
 \end{aligned}
\end{equation}
On the free abelianization it acts as
\begin{equation}\label{eq:gl-homology}
 a\longmapsto a,\qquad b\longmapsto b+d_0,\qquad
 c\longmapsto c-a,\qquad d_0\longmapsto d_0.
\end{equation}
\end{lemma}
\begin{proof}
All substitutions act on the rightmost expression first. Set
$D=r^{-1}d_0r$ and $s=D^{-1}rD$. From \eqref{eq:gl-inverse},
\[
 G(b)=D^{-1}b,\qquad G(r)=s,\qquad G(z)=r,\qquad
 G(c)=sr^{-1}cr,\qquad G(d_0)=D.
\]
For example,
$G(r)=(D^{-1}b)a(b^{-1}D)=D^{-1}rD$, and
$G(z)=D(D^{-1}rD)D^{-1}=r$. Substitution in
\eqref{eq:gl-forward} now gives
\[
 G(F(a))=a,\quad G(F(b))=DD^{-1}b=b,\quad
 G(F(c))=rs^{-1}(sr^{-1}cr)r^{-1}=c,
\]
and $G(F(d_0))=rDr^{-1}=d_0$.
Conversely, $F(r)=z$, $F(D)=z^{-1}(zd_0z^{-1})z=d_0$, and
$F(s)=d_0^{-1}zd_0=r$. Applying $F$ to the preceding expressions
for $G(b),G(c),G(d_0)$ gives $b,c,d_0$, respectively. Thus both
compositions are the identity on all four generators.

The boundary is preserved as a word, not just up to conjugacy:
\[
 \begin{aligned}
 F([a,b])&=a d_0r^{-1}d_0^{-1},\\
 F([c,d_0])&=zr^{-1}cd_0c^{-1}rd_0^{-1}z^{-1}.
 \end{aligned}
\]
Using $z=d_0rd_0^{-1}$, their product reduces to
$ar^{-1}cd_0c^{-1}d_0^{-1}=[a,b][c,d_0]$.
Finally $r$ and $z$ both have abelian image $a$, which gives
\eqref{eq:gl-homology} directly.
\end{proof}

\begin{proposition}[Descent to the full Galois group]
\label{prop:gl-full-descent}
Suppose $n\ge4$ is even. Apply Lemma~\ref{lem:gl-word} to
$(a,b,c,d_0)=(x_1,x_2,x_3,x_4)$ and fix all other generators,
including the tame factor. This defines a continuous automorphism
of $G_E$ with the action \eqref{eq:gl-homology} on abelianization.
More generally, a free-word automorphism of $x_1,\ldots,x_n$
descends in this manner if it fixes $x_1$ and the ordered
commutator product in \eqref{eq:gl-jw} exactly.
\end{proposition}
\begin{proof}
The substitution and its finite-word inverse extend continuously to
the free profinite product $V$ by its universal property. Both leave
the projection to $\mathfrak T$ unchanged, and hence preserve $Z$.
The kernel $J$ of the canonical maximal pro-$p$ quotient of $Z$
is characteristic in $Z$. Thus both substitutions preserve $J$ and
give inverse automorphisms of $\mathcal F$. This verifies the
wild-normal-subgroup condition in the original presentation.

They fix $\sigma,\theta,x_0$, and therefore fix $W_0$ and all its
profinite powers. They fix $x_1^{p^s}$ literally. The first two
commutators are fixed as their ordered product, and the later
commutators are fixed term by term. Thus the additional relator
and its closed normal closure are preserved. The induced maps on
the full quotient are inverse; the identification of that quotient
with $G_E$ is Jannsen--Wingberg's Theorem~2 with the algebraic closure
as extension. The same argument proves the general assertion.
\end{proof}

For clarity, the descended cross-handle automorphism is not inner.
The character to $\Z_p$ sending $x_4$ to $1$ and the other generators
and tame factor to $0$ satisfies the relative presentation. It has
value $0$ on $x_2$ and value $1$ on $F(x_2)=x_4x_2$; an inner
automorphism would preserve every abelian character. This observation
does not require any choice of logarithmic basis.

\subsection{The integral JW basis and the available linear actions}

We keep the same full presentation and its cyclotomic parameters when
passing to local reciprocity. Let $y_i$ be the logarithm of the
principal-unit image of $x_i$. Hoshi--Nishio
\cite[Lemma~1.3 and its proof]{HoshiNishio} prove that
$y_1,\ldots,y_n$ are a $\Q_p$-basis of $E$ and that
$y_0,\ldots,y_n$ generate $L_E$ over $\Z_p$. The former assertion
alone does not give an integral basis.

\begin{lemma}[Integral upgrade in the small tame range]
\label{lem:gl-integral-basis}
Assume $2\le e(E/\Q_p)\le p-2$. Then
$L_E=\mathfrak m_E$, and
\begin{equation}\label{eq:gl-integral-basis}
 v_i=y_i/p\quad(1\le i\le n)
\end{equation}
form a $\Z_p$-basis of $I_E$.
For even $n$, in these same coordinates,
\begin{equation}\label{eq:gl-trace-kernel}
 I_E\cap\ker\operatorname{Tr}_{E/\Q_p}
 =\Z_pv_1\oplus\bigoplus_{i=3}^n\Z_pv_i.
\end{equation}
If additionally $\operatorname{Tr} I_E=\Z_p$, then
$\operatorname{Tr}(v_2)$ is a unit.
\end{lemma}
\begin{proof}
The inequality $1/e>1/(p-1)$ makes logarithm and exponential inverse
on $1+\mathfrak m_E$ and $\mathfrak m_E$. The prime-to-$p$ torsion
units have zero logarithm, proving $L_E=\mathfrak m_E$.

Let $g,h$ be the lifts of the tame cyclotomic actions appearing in
the original presentation. With the notation of
\cite[Proposition~1.1]{HoshiNishio}, the mixed factor can be written
\[
 W_0=
 \left(x_0^{h^{p-1}}\theta x_0^{h^{p-2}}\theta
             \cdots x_0^h\theta\right)^{\epsilon_p g/(p-1)},
\]
where $\epsilon_p\in\widehat\Z$ is the projector onto $\Z_p$.
The tame relation makes the image of $\theta$ torsion in
abelianization, so it contributes zero to logarithm. Taking
logarithms in \eqref{eq:gl-jw} gives
\begin{equation}\label{eq:gl-y0}
 \left(1-\frac g{p-1}\sum_{r=1}^{p-1}h^r\right)y_0=p^s y_1.
\end{equation}
The reduction of $h$ modulo $p$ is not $1$. Indeed, trivial inertia
action on $\mu_p$ would make $E(\mu_p)/E$ unramified, forcing
$p-1$ to divide $e(E/\Q_p)$, contrary to the assumed range.
Thus the geometric sum in \eqref{eq:gl-y0} vanishes modulo $p$.
Its parenthesized coefficient is a unit, and
$y_0\in\Z_p y_1$. The remaining $y_i$ therefore generate $L_E$
integrally. Their known rational independence gives an integral
basis and proves \eqref{eq:gl-integral-basis}.

For even degree, Kondo's trace-kernel theorem
\cite[Theorem~1.3]{KondoGalois} identifies
$\ker\operatorname{Tr}$ with the rational span of
$v_1,v_3,\ldots,v_n$ in this JW basis. Intersecting with the proved
integral basis gives \eqref{eq:gl-trace-kernel}. Consequently the
trace ideal of $I_E$ is generated by $\operatorname{Tr}(v_2)$,
which proves the final assertion.
\end{proof}

Write $a=v_1$, $b=v_2$, and
$W=\bigoplus_{i=3}^n\Z_pv_i$. Let $\omega$ be the standard integral
alternating form on these remaining handle coordinates, with
$\omega(v_3,v_4)=1$. This form records the handle action; it is not
being identified with field multiplication. The covariant logarithmic
action of Proposition~\ref{prop:gl-full-descent} is
\[
 N_{v_4}(a)=a,\qquad N_{v_4}(b)=b+v_4,\qquad
 N_{v_4}(x)=x+\omega(v_4,x)a\quad(x\in W).
\]
The twist $x_2\mapsto x_2x_1$ gives $C_1(b)=b+a$, fixing $a,W$.

The boundary-preserving surface construction in
\cite[pp.~19--20 and Theorem~2.17]{KondoGalois} also gives an
individual full Galois lift of every integral symplectic matrix on
the remaining handles, fixing $\sigma,\theta,x_0,x_1,x_2$.
The construction on p.~19 explicitly includes even degree; the
odd-degree restriction later on p.~20 is not used here.
In particular, handle exchanges and rotations move $v_4$ to any
displayed basis direction. Conjugating the cross-handle lift
therefore realizes $N_w$ for every such direction.

For a general integer coordinate vector $w\in\Span_\Z\{v_3,\ldots,v_n\}$,
the linear formulas are
\begin{equation}\label{eq:gl-cross-linear}
 \begin{aligned}
 N_w(Aa+Bb+x)&=(A+\omega(w,x))a+Bb+x+Bw,\\
 C_c(Aa+Bb+x)&=(A+cB)a+Bb+x.
 \end{aligned}
\end{equation}
The following correction makes their realization precise.

\begin{lemma}[Ordered central correction]\label{lem:gl-central}
As identities of the induced linear maps,
\begin{equation}\label{eq:gl-central-law}
 N_wN_v=C_{\omega(w,v)}N_{w+v},\qquad
 C_cC_d=C_{c+d},\qquad C_cN_w=N_wC_c.
\end{equation}
For $w=\sum_i n_i e_i$ in the fixed ordered integer handle basis,
\begin{equation}\label{eq:gl-central-correction}
 N_w=C_{-\sum_{i<j}n_i n_j\omega(e_i,e_j)}
                N_{e_1}^{n_1}\cdots N_{e_{n-2}}^{n_{n-2}}.
\end{equation}
Thus every such $N_w$ is the canonical action of an actual full
Galois automorphism.
\end{lemma}
\begin{proof}
In coordinates $(A,B,x)$, applying $N_v$ and then $N_w$ gives
\[
 (A+\omega(v,x)+\omega(w,x+Bv),\ B,\ x+Bv+Bw).
\]
Bilinearity gives \eqref{eq:gl-central-law}; the other two identities
follow directly from \eqref{eq:gl-cross-linear}. Alternation implies
$N_w^{-1}=N_{-w}$ and $N_w^m=N_{mw}$ for every integer $m$.
An induction in the displayed left-to-right basis order gives
\eqref{eq:gl-central-correction}. Each basis cross and $C_1$ already
has a full lift. Integer powers, inverses and compositions of those
lifts realize the displayed linear image.
\end{proof}

The central identities are not asserted in the full nonabelian
automorphism group: lifts may differ by elements acting trivially
on abelianization. Nor is a group-homomorphic section of a mapping
class group into an outer Galois group required. Individual lifts
and the functoriality of their linear action suffice. In particular,
any prescribed vector modulo $p$ can be realized by choosing an
integer coordinate lift in \eqref{eq:gl-central-correction}.

\subsection{Covariance and integral Kummer transport}

For $\alpha\in\operatorname{Aut}(G_E)$, let $M_\alpha$ be the
canonical covariant action on principal-unit abelianization, extended
by logarithm and $\Q_p$-linearity to $E$. It sends each $y_i$ to
the logarithm of the image of $\alpha(x_i)$ and preserves $I_E$.
The norm/trace compatibility of this action is the one in
\cite[Lemma~2.3]{HoshiNishio}.

\begin{proposition}[The inverse in Kummer transport]
\label{prop:gl-variance}
For a continuous automorphism $\alpha$ and an integral equivariant
Tate-module isomorphism, let $F_\alpha$ be the induced logarithmic
Kummer pullback. Then
\begin{equation}\label{eq:gl-variance}
 F_\alpha=c_\alpha M_\alpha^{-1}
 \qquad\text{for some }c_\alpha\in\Z_p^\times.
\end{equation}
Thus the native valuations attained by $M_\alpha$ are attained by
one integral Kummer arrow belonging to $\alpha^{-1}$.
\end{proposition}
\begin{proof}
The cyclotomic character is determined by the local absolute Galois
group, so an equivariant integral Tate-module isomorphism exists;
its choices differ by $\Z_p^\times$
\cite[Proposition~1.1]{MochLocal}. Local duality identifies
$H^2(G_E,\Z_p(1))$ with $\Z_p$, and integral transport acts on it
by a unit $c_\alpha$. The cup-product reciprocity pairing gives,
for a continuous integral character $\chi$ and a Kummer class $z$,
\begin{equation}\label{eq:gl-cup-variance}
 \langle\chi\circ\alpha^{\rm ab},F_\alpha z\rangle
       =c_\alpha\langle\chi,z\rangle.
\end{equation}
Here the pairing is evaluation of $\chi$ on the pro-$p$ completion
of $E^\times$ under local reciprocity; see the conventions in
Proposition~\ref{prop:trace-transport} and
\cite[Chapter~III, Proposition~3.6 and \S4]{MilneCFT}.
The unit submodule is the annihilator of the unramified character,
so it is preserved. Torsion classes disappear after logarithm.

On the free unit lattice, the pullback character is $\chi\circ
M_\alpha$. These characters separate the lattice after rational
extension. Equation~\eqref{eq:gl-cup-variance} therefore says
$M_\alpha F_\alpha=c_\alpha\operatorname{id}$, proving
\eqref{eq:gl-variance}. Applying this to $\alpha^{-1}$ gives
$F_{\alpha^{-1}}=c_{\alpha^{-1}}M_\alpha$.
A $\Z_p$-unit scalar preserves the native valuation of every
vector, simultaneously. Uniform division of logarithmic coordinates
by $p$ does not alter the identity.
\end{proof}

\subsection{A degree-ten example}

\begin{lemma}[One minimal vector]\label{lem:gl-two-vectors}
Use the integral basis in
Lemma~\ref{lem:gl-integral-basis}, with even $n\ge4$, and let $\pi$ be
a uniformizer. Suppose
\[
 \operatorname{Tr}I_E=\operatorname{Tr}(\pi I_E)=\Z_p,
 \qquad u_0\in\pi I_E,\quad
 \operatorname{Tr}(u_0)\in\Z_p^\times,
\]
and $t_0\in I_E\setminus\pi I_E$ with
$\operatorname{Tr}(t_0)\in p\Z_p$. Then one full Galois
automorphism has canonical action $M$ satisfying
$Mu_0,Mt_0\in I_E\setminus\pi I_E$.
\end{lemma}
\begin{proof}
Put $\beta=\operatorname{Tr}(b)\in\Z_p^\times$. The respective
$b$-coordinates are $B_u=\operatorname{Tr}(u_0)/\beta\in\Z_p^\times$
and $B_t=\operatorname{Tr}(t_0)/\beta\in p\Z_p$.
If $a\notin\pi I_E$, use $C_1$. Modulo $\pi I_E$, its image of
$u_0$ is $B_ua\ne0$, and its image of $t_0$ is unchanged, because
$pI_E\subseteq\pi I_E$.

If $a\in\pi I_E$, then $I_E\cap\ker\operatorname{Tr}$ surjects
onto $I_E/\pi I_E$: subtract an element of $\pi I_E$ with the same
trace from any representative. Equation~\eqref{eq:gl-trace-kernel}
implies that some displayed basis vector $w\in W$ is outside
$\pi I_E$. The actual cross lift $N_w$ changes $u_0$ by
$B_uw+A_ua$ and $t_0$ by $B_tw+A_ta$, for integral $A_u,A_t$.
Modulo $\pi I_E$, the first image is $B_uw\ne0$, and the second
is the original nonzero image of $t_0$. The same map works for both.
\end{proof}

\begin{theorem}[A common minimum layer over $\Q_{19}$]
\label{thm:gl-nineteen}
Let
\[
 K_0=\Q_{19}(\mu_5),\qquad \pi^5=19,\qquad E=K_0(\pi),
\]
and put
\[
 u=\frac{\log20}{19},\qquad
 \tau=\frac{\log(1+19\pi)}{19},\qquad t=\tau/19.
\]
Then some full Galois automorphism has canonical action $M$ with
\begin{equation}\label{eq:gl-nineteen-min}
 v_{19}(Mu)=v_{19}(Mt)=-4/5.
\end{equation}
One integral Kummer arrow attains the same two valuations.
\end{theorem}
\begin{proof}
The order of $19$ modulo $5$ is $2$, so $K_0/\Q_{19}$ is
unramified quadratic. The Eisenstein polynomial $X^5-19$ and
$\mu_5\subset K_0$ give a Galois extension with $(n,e,f)=(10,5,2)$.
The logarithm range yields
\[
 I_E=\pi^{-4}\OO_E
     =\bigoplus_{r=-4}^{0}\OO_{K_0}\pi^r.
\]
The negative powers in this direct sum have relative trace zero,
and $\operatorname{Tr}_{E/K_0}(1)=5$. Hence
$\operatorname{Tr}I_E=\Z_{19}$. Also
$1/10\in\OO_E\subset\pi I_E$ has absolute trace $1$, proving
$\operatorname{Tr}(\pi I_E)=\Z_{19}$.

Leading logarithm terms give
$v_{19}(u)=0$ and $v_{19}(t)=-4/5$. The norm identity
\[
 \prod_{\zeta\in\mu_5}(1+19\zeta\pi)=1+19^6
\]
gives the exact traces
\[
 \operatorname{Tr}(u)=10u\in\Z_{19}^\times,
 \qquad \operatorname{Tr}(t)=\frac2{19^2}\log(1+19^6)
                         \in19^4\Z_{19}^\times.
\]
Lemma~\ref{lem:gl-two-vectors} proves \eqref{eq:gl-nineteen-min},
and Proposition~\ref{prop:gl-variance} gives the Kummer assertion.
\end{proof}

The covariant action in Theorem~\ref{thm:gl-nineteen} is not induced
by a field automorphism: it changes the native valuation of $u$
from $0$ to $-4/5$. This is an assertion about a specified local
vector, rather than a classification of all admissible linear maps.

\subsection{A simultaneous finite-family lemma}

The degree-ten argument begins with $t$ already at the minimum
layer. The square-index example below does not have this feature.
We therefore need a common operation on several vectors initially
lying above the minimum layer.

\begin{lemma}[Simultaneous parabolic reachability]
\label{lem:gl-finite-family}
Let $k$ be a finite field with $q$ elements, and let the finite-dimensional
$k$-vector space
\[
 V=ka\oplus kb\oplus W
\]
carry the standard nondegenerate alternating form on $W$, and let
$B$ be the $b$-coordinate. Let $L\subset V$ have codimension two,
and suppose $\ker B\to V/L$ is onto. Assume
\[
 U,T_1,\ldots,T_r\in L,\qquad B(U)\ne0,\qquad B(T_j)=0,
 \qquad T_j\notin ka,
\]
with $r+1<q$. The linear operations $C_c,N_w$ of
\eqref{eq:gl-cross-linear}, together with symplectic operations
on $W$ fixing $a,b$, contain one map sending
$U,T_1,\ldots,T_r$ all outside $L$.
If $k=\mathbf F_p$ and the coordinates come from the preceding
integral JW basis, one such map is induced by a full Galois
automorphism, using integer coordinate lifts.
\end{lemma}
\begin{proof}
Write $T_j=A_ja+w_j$. The hypotheses imply $w_j\ne0$, so each
functional $w\mapsto\omega(w,w_j)$ is nonzero.

Suppose first that $a\notin L$. Fewer than $q$ proper hyperplanes
cannot cover $W$: in dimension $D$, their union has fewer than
$q\cdot q^{D-1}=q^D$ elements. Choose $w$ with
$\omega(w,w_j)\ne0$ for all $j$. Then
\[
 N_w(T_j)+L=\omega(w,w_j)(a+L)\ne0.
\]
If $N_w(U)\notin L$, the construction is complete. Otherwise
compose with $C_1$. It adds the nonzero class $B(U)(a+L)$ to
$N_w(U)$ and leaves all the $T_j$ images unchanged because their
$B$-coordinates remain zero.

Suppose now that $a\in L$. Then $\lambda:W\to V/L$ is onto.
Also $\lambda(w_j)=0$ for every $j$. Avoid $\ker\lambda$ and
the $r$ hyperplanes $\ker\omega(-,w_j)$ simultaneously, obtaining
$s\in W$ such that
\[
 \lambda(s)\ne0,\qquad\omega(s,w_j)\ne0\quad\text{for all }j.
\]
There are $r+1<q$ proper subspaces, so the same union bound applies.
The single symplectic transvection
\[
 S_s(x)=x+\omega(s,x)s\quad(x\in W),\qquad
 S_s(a)=a,\quad S_s(b)=b
\]
has inverse given by changing the plus sign to a minus sign.
Expanding its pairing shows it preserves $\omega$; the cross
terms cancel and $\omega(s,s)=0$. Since the original projections
vanish,
\[
 S_s(T_j)+L=\omega(s,w_j)\lambda(s)\ne0
\]
for every $j$. If $S_s(U)\notin L$, stop. Otherwise choose a
displayed basis vector $e$ of $W$ with $\lambda(e)\ne0$ and
apply $N_e$. Its change on the zero projection of $S_s(U)$ is
$B(U)\lambda(e)\ne0$. It does not change any $T_j$ projection,
because $B(T_j)=0$ and $a\in L$.

For the arithmetic assertion, lift the chosen residue vectors to
integer coordinate vectors. Lemma~\ref{lem:gl-central} realizes
each required cross. The lifted transvection is an element of
$\operatorname{Sp}(W_\Z)$ and has an individual full Galois lift
by the surface construction cited above. Compose these lifts.
The argument uses one map for the whole family, not a different
map for each vector.
\end{proof}

\subsection{Three square labels at the prime 139}

We state the next example for a general unit part of the parameter.
Let $p=139$ and choose any $b_0\in\Q_p$ with $v_p(b_0)=1$. Put
\begin{equation}\label{eq:gl-139-field}
 K_0=\Q_p(\mu_{210}),\qquad \pi^{105}=b_0,
 \qquad E=K_0(\pi),\qquad q_0=b_0^2.
\end{equation}
Thus $\pi^{210}=q_0$ and $\pi^{15}=q_0^{1/14}$ with these choices
of roots. The field in \eqref{eq:gl-139-field} is exactly
$\Q_p(\mu_{210},q_0^{1/210})$; no global curve with a prescribed
exact parameter is assumed here.

\begin{theorem}[One full Galois arrow for three square labels]
\label{thm:gl-three-labels}
In \eqref{eq:gl-139-field}, set
\begin{equation}\label{eq:gl-139-vectors}
 \begin{aligned}
 u&=\frac{\log(1+p)}p, &
 \tau_s&=\frac{\log(1+p\pi^{15s})}p,\\
 k_s&=\left\lfloor\frac s7+\frac{104}{105}\right\rfloor,&
 T_s&=p^{-k_s}\tau_s\qquad(s=1,4,9).
 \end{aligned}
\end{equation}
Then $k_1=k_4=1$, $k_9=2$. One full Galois automorphism has
canonical action $M$ such that
\begin{equation}\label{eq:gl-139-minimum}
 v_p(Mu)=v_p(MT_1)=v_p(MT_4)=v_p(MT_9)=-104/105.
\end{equation}
Consequently
\begin{equation}\label{eq:gl-139-theta}
 v_p(M\tau_1)=1/105,\qquad v_p(M\tau_4)=1/105,
 \qquad v_p(M\tau_9)=106/105.
\end{equation}
One integral Kummer arrow attains all these values simultaneously.
\end{theorem}
\begin{proof}
Since $139^2\equiv1\pmod{210}$ and $139\not\equiv1\pmod{210}$,
$K_0/\Q_p$ is unramified quadratic. The Eisenstein polynomial
$X^{105}-b_0$ and $\mu_{105}\subset K_0$ give a Galois field
with $(n,e,f)=(210,105,2)$. As $1/105>1/138$,
\begin{equation}\label{eq:gl-139-lattice}
 I_E=\pi^{-104}\OO_E,
 \qquad pI_E=\pi\OO_E,
 \qquad \pi I_E=\pi^{-103}\OO_E.
\end{equation}
These are equalities of fractional ideals: $b_0/p$ is a rational
unit and need not be $1$.

The power basis gives
$I_E=\bigoplus_{r=-104}^0\OO_{K_0}\pi^r$.
All negative powers in this sum have relative trace zero and
$\operatorname{Tr}_{E/K_0}(1)=105$. Hence
$\operatorname{Tr}I_E=\Z_p$. Also $1/210\in\OO_E\subset\pi I_E$
has trace $1$. Subtracting $\operatorname{Tr}(z)/210$ from
$z\in I_E$ proves that $I_E\cap\ker\operatorname{Tr}$ surjects
onto $I_E/\pi I_E$. Lemma~\ref{lem:gl-integral-basis} applies.

The integer $k_s$ is exactly the largest $k$ for which
$\tau_s\in p^k I_E$. Leading logarithm terms and the exact norm
give the following data:
\[
 \begin{array}{c|ccc}
 s&k_s&v_p(T_s)&v_p(\operatorname{Tr}T_s)\\ \hline
 1&1&-90/105&6\\
 4&1&-45/105&9\\
 9&2&-75/105&13
 \end{array}
\]
Indeed, for these $s$, the seven conjugates of $\pi^{15s}$ occur
fifteen times in the relative norm, and therefore
\begin{equation}\label{eq:gl-139-trace}
 \begin{aligned}
 \operatorname{Nm}_{E/K_0}(1+p\pi^{15s})&=(1+p^7b_0^s)^{15},\\
 \operatorname{Tr}(T_s)
     &=\frac{30}{p^{k_s+1}}\log(1+p^7b_0^s).
 \end{aligned}
\end{equation}
The last expression has valuation $s+6-k_s$, as displayed.
Furthermore $\operatorname{Tr}(u)=210\log(1+p)/p$ is a unit.

Pass to the $\mathbf F_p$-space
\[
 V=I_E/pI_E,
 \qquad L=\pi I_E/pI_E,
 \qquad \dim V=210,\quad\operatorname{codim}L=2.
\]
The integral JW basis gives $V=\mathbf F_pa\oplus\mathbf F_pb
\oplus W$, with $\dim W=208$ and $B$ its $b$-coordinate.
Equation~\eqref{eq:gl-trace-kernel} and
$\operatorname{Tr}(b)\in\Z_p^\times$ imply
\[
 B(u)\ne0,\qquad B(T_s)=0,
 \qquad u,T_s\in L,
 \qquad T_s\ne0\text{ in }V.
\]
The trace-zero subspace surjects onto $V/L$ by the subtraction
argument above. There remains a possible fixed-line obstruction:
one or more of the $T_s$ might belong to $\mathbf F_pa$.

We remove it by one common field automorphism. The logarithm tail
satisfies
\begin{equation}\label{eq:gl-139-leading}
 T_s\equiv\pi^{15s}/p^{k_s}\pmod{pI_E}.
\end{equation}
For its term of degree $m\ge2$, use $v_p(m)\le m-2$ to obtain
\[
 m(1+s/7)-(k_s+1)-v_p(m)
 \ge1+ms/7-k_s\ge1+2s/7-k_s.
\]
The last three bounds are $2/7,8/7,11/7$, all larger than
$1/105$, the minimum valuation of $pI_E$.

Let $\sigma(\pi)=\zeta_{105}\pi$, fixing $K_0$. For
$h=0,\ldots,6$, equation~\eqref{eq:gl-139-leading} gives
\[
 \sigma^h(T_s)\equiv
 \zeta_7^{sh}\pi^{15s}/p^{k_s}\pmod{pI_E}.
\]
For each $s$, these determine seven distinct $\mathbf F_{139}$-lines.
To see this, the reduction of $\zeta_7$ has order seven and
$7\nmid138$, so its group meets $\mathbf F_{139}^\times$ only in
$1$. Equality of two lines would force their leading coefficients
to have a ratio in that intersection. A nonzero residue difference
would be a unit times $\pi^{15s}/p^{k_s}$, whose negative valuation
precludes membership in $pI_E$. Since $s$ is prime to seven, the
two exponents must be equal modulo seven.

For each of the three labels, at most one of these seven values of
$h$ puts $\sigma^h(T_s)$ in the fixed line $\mathbf F_pa$.
Avoid the at most three excluded values. This one $\sigma^h$
fixes $u$, preserves $L$ and trace, and places all three labels
outside the fixed line. It is an actual field automorphism, hence
has a representative on the full local Galois group.

Apply Lemma~\ref{lem:gl-finite-family} to these three resulting
vectors and $u$. The numerical condition is $3+1<139$.
More precisely, in its second case the excluded sets have at most
$139^{206}+3\cdot139^{207}<139^{208}$ elements. The lemma gives
one further actual lift sending all four vectors outside $L$.
Composing with the chosen inertia representative proves
\eqref{eq:gl-139-minimum}. Linearity and $\tau_s=p^{k_s}T_s$
give \eqref{eq:gl-139-theta}; Proposition~\ref{prop:gl-variance}
supplies the common Kummer arrow.
\end{proof}

The unit $b_0/p$ is fixed by inertia. Thus the proof includes all
parameters in \eqref{eq:gl-139-field}, rather than only
$b_0=p$. The separate passage from a rational Frey curve to these
local data does not enter the simultaneous-reachability argument.

\subsection{Exact native maximal-order hulls}

Let $E/\Q_p$ now be either Galois field above. For $m\ge1$, write
\[
 \mathcal T_m=E^{\otimes_{\Q_p}m},\qquad
 B_m=\prod_{\operatorname{Gal}(E/\Q_p)^{m-1}}\OO_E,
\]
using the decomposition
\begin{equation}\label{eq:gl-tensor-decomposition}
 x_1\otimes\cdots\otimes x_m\longmapsto
 \bigl(x_1\gamma_2(x_2)\cdots\gamma_m(x_m)\bigr)_
                {(\gamma_2,\ldots,\gamma_m)}.
\end{equation}
Thus $B_m$ is the maximal order in this fixed finite etale carrier.
In $\pi^rB_m$, the uniformizer is embedded through the first tensor
factor, so it has the same valuation in every component.
Let $\Gamma_E$ be the set of integral logarithmic Kummer transports
arising from automorphisms of the same $G_E$, with integral
Tate-module identifications. Every $F\in\Gamma_E$ is $\Q_p$-linear
and preserves $I_E$ and $p^kI_E$ for all integers $k$.

\begin{proposition}[An attained block hull]\label{prop:gl-block-hull}
Suppose $I_E=\pi^{1-e}\OO_E$, and put $\kappa=1-1/e$.
Let $u,t\in I_E$, $\tau=p^kt$ with integer $k$, and suppose one
$F_0\in\Gamma_E$ satisfies $v_p(F_0u)=v_p(F_0t)=-\kappa$.
Define the coherent repeated-label family
\[
 S_m=\{F(\tau)\otimes F(u)\otimes\cdots\otimes F(u):
                      F\in\Gamma_E\},\qquad
 H_m=\Span_{B_m}(S_m).
\]
Then
\begin{equation}\label{eq:gl-block-hull}
 H_m=\pi^{ek-(e-1)m}B_m.
\end{equation}
The equality also holds after taking topological closure. More
generally, it remains true if the set of source tuples is enlarged
to any set lying between the displayed coherent tuples and
$p^kI_E\times I_E^{m-1}$.
\end{proposition}
\begin{proof}
For every $F$, the first factor has valuation at least $k-\kappa$,
and each remaining factor at least $-\kappa$. All field embeddings
in \eqref{eq:gl-tensor-decomposition} preserve valuation. Every
component of every tensor generator therefore has valuation at least
$k-m\kappa$, proving containment in the right-hand side of
\eqref{eq:gl-block-hull}.

The single arrow $F_0$ gives equality in this valuation in every
component. Its tensor is consequently
$\pi^{ek-(e-1)m}$ times a unit of the product ring $B_m$.
Its principal $B_m$-span is the entire right-hand side, giving
the reverse containment. This fractional product lattice is
closed. The same upper bound holds for every tuple in the stated
containing product, while the original attaining tuple remains
present; this proves the enlargement assertion.
\end{proof}

The enlargement includes the closed $\Z_p$-convex hull taken in
the product of the source coordinate spaces, since that product
lattice is closed and $\Z_p$-convex. This does not identify such
a convex hull with a $B_m$-module hull: tensor formation and the
final $B_m$-span remain distinct operations.

\begin{corollary}[The two exact local calculations]
\label{cor:gl-exact-hulls}
For Theorem~\ref{thm:gl-nineteen}, the family with one $\tau$
factor and $m-1$ background factors has
\[
 H_m=\pi^{5-4m}B_m\quad(m\ge1),
 \qquad\text{in particular }H_3=\pi^{-7}B_3.
\]
For Theorem~\ref{thm:gl-three-labels}, let $H_j$ use one
$\tau_{j^2}$ factor and $j$ background factors, for $j=1,2,3$.
Then one and the same arrow attains all three hulls, and
\begin{equation}\label{eq:gl-three-hulls}
 H_1=\pi^{-103}B_2,\qquad
 H_2=\pi^{-207}B_3,\qquad
 H_3=\pi^{-206}B_4.
\end{equation}
\end{corollary}
\begin{proof}
Use Proposition~\ref{prop:gl-block-hull} with $(e,k)=(5,1)$
in the first case. In the second case $e=105$, $m=j+1$, and
$(k_1,k_4,k_9)=(1,1,2)$. Thus the three exponents are
$105k_{j^2}-104(j+1)=-103,-207,-206$.
The common witness is supplied by Theorem~\ref{thm:gl-three-labels}
and the single inverse/unit conversion of
Proposition~\ref{prop:gl-variance}.
\end{proof}

These nonintegral hulls are exact calculations for stated local
source families. They do not establish a comparison between different
native markings, a cross-Frobenius identification, or compatible
global initial theta data. In particular, identifying the source
sets needed to compare a published pilot container with one of these
hulls remains a separate requirement.

The free-word identities, the discrete free-group automorphism,
the linear central correction, and the finite-family transvection
statements have independent Lean proofs. The relative profinite
quotient, local class field theory, local duality, integral JW basis,
and logarithmic norm calculations are mathematical source arguments,
not Lean declarations of local-field theorems in this work. The
global abc bound remains an additional obligation.
